\chapter[Efficient Question Answering over External Knowledge][Efficient \acs{qa} over External Knowledge]{Efficient Question Answering over External Knowledge}
\label{chap:refactx}


As described in \Cref{sec:intro,sec:rw:llms}, \acp{llm} have demonstrated a variety of capabilities, ranging from natural language understanding and generation to reasoning, especially when enhanced with techniques like \acf{cot} prompting~\cite{wei2022chainofthought}. However, they are prone to hallucinations~\cite{maynez-etal-2020-faithfulness-hallucination} and their internal knowledge remains limited to their training data (see \Cref{sec:llms_vs_kbs}). This complicates their application to knowledge-intensive tasks such as \acf{qa}, especially when these tasks necessitate accessing information beyond the parametric knowledge of \acp{llm}---for example, recent data or domain-specific data that is not publicly available.


Solutions to augment \acp{llm} with external knowledge exist, as reviewed in \Cref{sec:sota:qa}, where we categorized between \emph{input-based} and \emph{latent-interaction} knowledge injection techniques. However, the former rely on external retrievers or pipelines that increase latency, system complexity, and suffer from error propagation.

In the latter case, instead, models acquire external knowledge more internally, avoiding the need for external retrievers, but they require architectural modifications complicating the use of recent \acp{llm}. 

\ac{qa} is relevant for knowledge-intensive domains and, thus, for this thesis work, and appears in the use cases we considered in \Cref{sec:intro}. However, regulations on private data could forbid the use of powerful \ac{api}-based \acp{llm}, forcing users to rely on executing models locally with hardware-related constraints. For this reason, the objective \ref{ob:refactx} aims at developing an efficient and scalable \ac{qa} system that additionally allows users to trace and verify the received answer to the original documents or pieces of knowledge.

Addressing \ref{ob:refactx}, this chapter describes Reliable Fact eXtractor (\refact), a \ac{qa} approach able to integrate external knowledge without depending on additional models or external retrievers, but relying solely on constrained generation, supported by a pre-built prefix-tree index, designed to facilitate and speed-up the acquisition of external facts.

A few recent approaches have applied \emph{constrained generation} to \ac{qa}, achieving promising results~\cite{li-etal-2025-decoding-graphs,luo2024graphconstrainedreasoningfaithfulreasoning}, however, they did not focus on the scalability to large knowledge bases.
This technique restricts the model's output space during decoding to sequences that satisfy predefined structural, syntactic, or semantic constraints. In our case, we use it for ensuring the \ac{llm} generates a fact that really exist in the \ac{kg}, with the aim of preventing hallucinations.

At inference time, with \refactx the \ac{llm} is instructed via \acf{icl}~\cite{dong2024survey} to invoke the \emph{Fact} command when external facts are required. Once the command is recognized, constrained generation is activated: the model produces tokens only along valid paths in the prefix tree, guaranteeing that the output matches a fact from the knowledge base, which is derived from Wikidata~\cite{wikidata}. Once an entire \ac{kb} fact is generated, the decoding mechanism reverts back to normal.
It is important to note that every generated token is selected based on \ac{llm}'s probability estimates. Constrained generation only narrows the vocabulary to the tokens that form an existing fact in the \ac{kb}, but always leaves the final choice to the \ac{llm}.


\begin{figure}[H]
    \centering
    \includegraphics[width=\textwidth]{images/refact_example.pdf}
    \caption{\refactx answering an open-domain question.
The \ac{llm} sketches a plan, makes two \emph{Fact} calls, and---with constrained generation (blue underline)---inserts valid facts from a Wikidata-based prefix tree before giving the final answer.
    }
    \label{fig:example}
    % \vspace{2em}
    \end{figure}


\paragraph{Motivating Example.} The approach is illustrated in \Cref{fig:example}, by depicting an example with the question \textit{``When was the director of Slumdog Millionaire born?''}. The model is instructed to first reason on how to reach the correct answer, and then it starts to acquire facts.
% TODO fact command?
Once the \emph{Fact} command is called, constrained generation is enabled (underlined in blue) and the model is guided to generate an existing fact. This way the model is able to first find that \textit{``Danny Boyle''} is the director of \textit{``Slumdog Millionaire''}, then to find his birthdate, and finally to answer the question correctly.

Furthermore, the facts generated can be highlighted, as in \Cref{fig:example}, so that users are aware that those facts exist in the \ac{kb}, supporting traceability and verifiability (\ref{ch:traceability}).

Although not enforced, the \acl{dia} design proposed in \Cref{chap:architecture} supports the use of \aclp{kb} that contain facts, which can be required by advanced systems such as \refactx.

% \newpage
The main contributions of this chapter can be summarized as follows:
\begin{itemize}
  \item \textbf{\refactx}: a generic, constrained-generation system that enables any \ac{llm} to access very large \acp{kb}, with no external retrievers or pipelines.
  \item \textbf{Efficiency and scalability}: a disk-backed prefix tree capable of holding up to 800 million facts, adding only $\sim$1\% latency.  
  \item \textbf{Empirical validation}: insights on four \ac{qa} datasets suggest that \refactx can achieve competitive results, with accuracy gains exceeding 20 percentage points at over 90\% precision on certain benchmarks---compared to \acp{llm} answering only with their parametric knowledge.
\end{itemize}

The following sections describe the decoding mechanism (\Cref{sec:constrained-gen}), the scalable Wikidata index (\Cref{sec:wikidata-repr}), and they are integrated into \acl{qa} workflow (\Cref{sec:refactx-qa}). Then, in we describe the experimental setup ( \Cref{sec:rx:experiments}), present (\Cref{sec:rx:results}) and discuss (\Cref{sec:rx:discussion}) the results. Finally, we conclude with \Cref{sec:rx:conclusion}. 

This chapter is related to the publication \textcite{Pozzi2025ReFactX}. The corresponding implementation is released under the Apache~2.0 license on GitHub\footnote{\url{https://github.com/rpo19/ReFactX}}.

